\chapter{نشانگر زنده‌بودن و آزمون دو چراغ}

\section{هدف آزمون}
دو چراغ متصل به پایه‌های
\lr{PC13}
و
\lr{PC14}
یک نشانگر مستقل و غیرمسدودکننده برای مشاهده اجرای برنامه فراهم می‌کنند. الگوی عادی فقط زمانی جلو می‌رود که تابع
\lr{App\_Process}
مرتب در حلقه اصلی فراخوانی شود. ثابت‌ماندن طولانی الگو، پس از درنظرگرفتن زمان طبیعی عملیات حسگرها، می‌تواند نشانه توقف حلقه اصلی یا گیرکردن یک تابع باشد. الگوی خطای سریع نیز ورود برنامه به
\lr{Error\_Handler}
را از اجرای عادی متمایز می‌کند.

این نشانگر جایگزین واچ‌داگ، ثبت خطا یا دیباگر نیست؛ چراغ فقط یک مشاهده سریع میدانی از زنده‌بودن حلقه اصلی ارائه می‌دهد.

\section{اتصال الکتریکی و قطب فعال}
هر چراغ باید مقاومت سری مستقل داشته باشد. برای کاهش بار پایه‌های دامنه پشتیبان درگاه سوم، مقاومت ۶۸۰ اهم تا یک کیلواهم و چراغ کم‌جریان پیشنهاد می‌شود. اتصال پیش‌فرض نرم‌افزار فعال‌پایین است:

\begin{center}
\lr{\texttt{3.3 V -> series resistor -> LED anode -> LED cathode -> PC13 or PC14}}
\end{center}

در این اتصال، صفر منطقی چراغ را روشن می‌کند. اگر اتصال برد از خروجی پایه به مقاومت و سپس چراغ و زمین است، مقدار
\lr{DIAGNOSTIC\_LED\_ACTIVE\_LOW}
در فایل
\filepath{Core/Inc/diagnostic\_led.h}
باید از یک به صفر تغییر کند. دو چراغ نباید یک مقاومت مشترک داشته باشند. زمین برد و پروگرامر نیز باید مشترک باشد.

\section{تنظیمات مولد کد}
در فایل
\filepath{STM32F107.ioc}
هر دو پایه به‌صورت خروجی پوش‌پول با سرعت پایین، بدون مقاومت کشش داخلی و سطح اولیه بالا ثبت شده‌اند. برچسب‌ها به‌ترتیب
\lr{DEBUG\_LED\_1}
و
\lr{DEBUG\_LED\_2}
هستند. منبع کلاک کم‌فرکانس خارجی فعال نیست؛ بنابراین استفاده از
\lr{PC14}
به‌عنوان خروجی با تنظیم کلاک موجود تعارض ندارد.

ماژول عیب‌یابی هنگام راه‌اندازی همین تنظیم را دوباره و به‌صورت بی‌خطر اعمال می‌کند. این تکرار عمدی است تا اگر خطای راه‌اندازی پیش از پایان تابع تنظیم پایه‌ها رخ داد، امکان نمایش خطا باقی بماند. پس از هر تولید مجدد باید وجود دو پایه، برچسب‌ها، سطح اولیه بالا و سرعت پایین کنترل شود.

\section{فایل‌ها و رابط عمومی}
پیاده‌سازی در
\filepath{Core/Src/diagnostic\_led.c}
و رابط عمومی در
\filepath{Core/Inc/diagnostic\_led.h}
قرار دارد. فایل پیاده‌سازی مستقل از فایل‌های مولد است و به پروژه
\lr{Keil}
نیز افزوده شده است؛
\lr{PlatformIO}
آن را به‌طور خودکار از پوشه سورس پیدا می‌کند.

\subsection{تابع \lr{DiagnosticLed\_Init}}
ورودی و خروجی ندارد. کلاک درگاه
\lr{GPIOC}
را فعال، دو پایه را خروجی پوش‌پول کم‌سرعت و هر دو چراغ را خاموش می‌کند. چند بار فراخوانی آن مجاز است، ولی هر بار حالت و گام الگو را بازنشانی می‌کند.

\subsection{تابع \lr{DiagnosticLed\_SetMode}}
یک مقدار از نوع
\lr{DiagnosticLed\_Mode}
می‌گیرد، الگوی جدید را انتخاب و اجرای آن را از گام نخست آغاز می‌کند. مقدار خارج از محدوده به حالت خاموش تبدیل می‌شود. این تابع برای متن اصلی طراحی شده و نباید هم‌زمان از وقفه تغییر داده شود.

\subsection{تابع \lr{DiagnosticLed\_Process}}
ورودی و خروجی ندارد و باید با بیشترین تکرار ممکن در حلقه اصلی صدا زده شود. زمان را با
\lr{HAL\_GetTick}
می‌خواند و با تفریق بدون علامت، عبور شمارنده میلی‌ثانیه از بیشینه را نیز درست مدیریت می‌کند. اگر موعد گام بعد نرسیده باشد بلافاصله بازمی‌گردد؛ در نتیجه هیچ تأخیر مسدودکننده‌ای به سامانه تحمیل نمی‌شود.

اگر حلقه اصلی برای چند گام متوقف شود، تابع در اولین اجرای بعدی تعداد گام‌های گذشته را محاسبه و الگو را به موقعیت درست می‌برد. این رفتار از انباشته‌شدن تأخیر زمانی جلوگیری می‌کند.

\subsection{تابع \lr{DiagnosticLed\_GetMode}}
ورودی ندارد و حالت فعال را برمی‌گرداند. این تابع برای مشاهده در دیباگر، گزارش وضعیت یا توسعه فرمان‌های عیب‌یابی آینده درنظر گرفته شده است و سخت‌افزار را تغییر نمی‌دهد.

\section{الگوهای قابل انتخاب}
\subsection{حالت \lr{DIAGNOSTIC\_LED\_MODE\_OFF}}
هر دو چراغ خاموش می‌شوند. این حالت برای غیرفعال‌کردن آزمون بدون حذف کد است.

\subsection{حالت \lr{DIAGNOSTIC\_LED\_MODE\_HEARTBEAT}}
دو تپ کوتاه و سپس مکث ایجاد می‌شود. این الگو نشانگر آرام برای بهره‌برداری دائمی است.

\subsection{حالت \lr{DIAGNOSTIC\_LED\_MODE\_DANCE}}
چراغ اول، چراغ دوم، هر دو و سپس مکث اجرا می‌شود. چرخه تقریباً یک ثانیه است و حالت پیش‌فرض برنامه محسوب می‌شود.

\subsection{حالت \lr{DIAGNOSTIC\_LED\_MODE\_ERROR}}
دو چراغ هر ۱۵۰ میلی‌ثانیه به‌سرعت جابه‌جا می‌شوند. این الگو نشان‌دهنده ورود به خطای بحرانی است.

\section{اتصال به چرخه برنامه}
تابع
\lr{App\_Init}
ماژول را راه‌اندازی و حالت رقص را انتخاب می‌کند. تابع
\lr{App\_Process}
پیش از رسیدگی به فرمان‌ها و نمونه‌برداری، یک بار تابع پردازش چراغ را اجرا می‌کند. اگر برنامه وارد
\lr{Error\_Handler}
شود، دو پایه دوباره راه‌اندازی و حالت خطا در یک حلقه اختصاصی پردازش می‌شود.

برای استفاده آرام‌تر پس از پایان آزمون کافی است در
\filepath{Core/Src/app.c}
حالت اولیه از رقص به ضربان تغییر کند؛ حذف فراخوانی پردازش توصیه نمی‌شود، زیرا ارزش نشانگر زنده‌بودن از بین می‌رود.

\section{پروگرام با رابط چهار سیمه}
برای پروگرام معمولی اتصال‌های زمین مشترک، ولتاژ مرجع ۳.۳ ولت، داده رابط و کلاک رابط لازم‌اند. نام دو خط رابط به‌ترتیب
\lr{SWDIO}
روی
\lr{PA13}
و
\lr{SWCLK}
روی
\lr{PA14}
است. در صورت دشواری اتصال، خط بازنشانی نیز می‌تواند به‌عنوان سیم پنجم افزوده شود.

پس از اتصال
\lr{J-Link}
و روشن‌کردن هدف، فرمان‌های زیر خروجی را ساخته و پروگرام می‌کنند:

\begin{LTR}
\begin{verbatim}
pio run
pio run --target upload
\end{verbatim}
\end{LTR}

در
\lr{Keil}
نیز هدف
\lr{STM32F107}
باز، پروژه ساخته و گزینه بارگذاری انتخاب می‌شود. فایل
\filepath{MDK-ARM/STM32F107.uvprojx}
از قبل ماژول چراغ را در گروه کد کاربر دارد.

\section{نتیجه پروگرام انجام‌شده}
در تاریخ ۱۴۰۵/۰۶/۲۳، بارگذاری مستقیم از همین پروژه با
\lr{PlatformIO}
و
\lr{J-Link}
انجام شد. پروگرامر از طریق رابط
\lr{USB}
شناسایی و ولتاژ مرجع هدف ۳.۳۲۵ ولت گزارش شد. قطعه
\lr{STM32F107VC}
و هسته
\lr{Cortex-M3}
از مسیر
\lr{SWD}
با موفقیت شناسایی شدند. پاک‌کردن، نوشتن، تطبیق محتوای حافظه و بازنشانی نهایی همگی بدون خطا پایان یافتند. بنابراین فایل اجرایی این نسخه روی سخت‌افزار بارگذاری شده است؛ مشاهده فیزیکی ترتیب چراغ‌ها باید توسط کاربر روی برد تأیید شود.

\section{روند آزمون و تفسیر نتیجه}
ابتدا برای جلوگیری از عملکرد ناخواسته، بارهای قدرت رله، هیتر و سرمایش قطع یا ایمن شوند. سپس برنامه بارگذاری و برد بازنشانی شود.

\begin{enumerate}
  \item در ثانیه نخست باید ترتیب چراغ اول، چراغ دوم و هر دو مشاهده شود.
  \item ادامه منظم چرخه نشان می‌دهد حلقه اصلی در حال اجرا و تیک سامانه فعال است.
  \item مکث یا پرش کوتاه می‌تواند ناشی از عملیات مسدودکننده فعلی حسگر دما یا ارسال سریال باشد و باید با دیباگر زمان‌سنجی شود.
  \item ثابت‌شدن طولانی الگو یعنی حلقه اصلی دیگر تابع پردازش را اجرا نمی‌کند یا تیک سامانه متوقف شده است.
  \item تناوب سریع و پیوسته چراغ‌ها یعنی برنامه وارد هندلر خطای بحرانی شده است.
  \item خاموش‌بودن دائمی هر دو چراغ می‌تواند از قطب اشتباه، سیم‌بندی نادرست، نبود مقاومت و تغذیه، پروگرام‌نشدن یا اجرا‌نشدن برنامه ناشی شود.
\end{enumerate}

این آزمون فقط راه‌اندازی اولیه و زنده‌بودن نرم‌افزار را نشان می‌دهد. تأیید سنسورها، ارتباط سریال، حافظه و کنترل بارها باید مطابق چک‌لیست آزمون سامانه جداگانه انجام شود.
